\documentclass[12pt]{article}

\usepackage[margin=1.2in]{geometry}
\usepackage{amsmath,amsthm,amssymb}
\usepackage{microtype}
\usepackage[authoryear,round]{natbib}
\usepackage{xcolor}
\definecolor{linkgreen}{RGB}{0,120,100}
\usepackage[colorlinks=true,linkcolor=linkgreen,citecolor=linkgreen,urlcolor=linkgreen]{hyperref}

% Reduce title font from \LARGE (default) to \Large so title fits in 2 lines
\makeatletter
\renewcommand{\maketitle}{%
  \begin{center}%
    {\Large\bfseries \@title \par}%
    \vskip 1.5em%
    {\normalsize \@author \par}%
    \vskip 0.8em%
    {\normalsize \@date}%
  \end{center}%
  \vskip 2em%
}
\makeatother

% Theorem environments
\theoremstyle{plain}
\newtheorem{proposition}{Proposition}
\newtheorem{lemma}{Lemma}
\newtheorem{corollary}{Corollary}
\theoremstyle{definition}
\newtheorem{assumption}{Assumption}
\newtheorem{definition}{Definition}
\theoremstyle{remark}
\newtheorem{remark}{Remark}

\title{Education Investment under China's Hukou System:\\
A Theoretical Framework}
\author{%
  \texttt{pAI-Econ-claude} (\texttt{theoretical-economics-claude-skill})\\[4pt]
  \small\url{https://github.com/maxwell2732/pAI-Econ-claude}\\[2pt]
  \small Claude Sonnet 4.6%
}
\date{Draft: June 14, 2026}

\begin{document}

\maketitle

\section{Introduction}

This paper develops a formal model of educational investment under China's hukou
(household registration) system.  A rural child and an urban child of identical
innate ability $\theta$ face different expected returns to education: the rural
child earns a fraction $\delta\in(0,1)$ of the urban wage unless she crosses an
elite education threshold $\hat{e}$, at which point she gains probability
$p_{0}>0$ of accessing the full urban labor market.  We embed this two-regime
payoff structure into the \citet{Becker1964} human capital framework and
characterize the resulting investment patterns.

Our central finding is that the hukou system generates a \emph{non-monotone}
relationship between ability and investment distortion.  Low-ability rural
children rationally under-invest relative to equally-able urban peers, confirming
the benchmark of \citet{Arrow1973}: a proportional wage penalty suppresses the
marginal return to education and reduces optimal investment.  Intermediate-ability
rural children, however, rationally \emph{over-invest}.  For these children,
whose unconstrained optimum lies just below $\hat{e}$, the escape-route
premium---the expected wage gain from crossing the elite threshold---outweighs
the cost of over-shooting.  The institutional barrier does not uniformly suppress
investment; for those near enough to the threshold to find the gamble worthwhile,
it reverses the direction of distortion.

Two further results follow.  First, the welfare loss from hukou is non-monotone
in ability, peaking among intermediate-ability children near $\hat{e}$ who bear
the cost of over-investment with certainty while receiving the escape premium only
with probability $p_{0}$.  Second, hukou reform that equalizes returns---as
China's post-2014 reforms aim to do---simultaneously improves welfare for all
rural children and reduces educational attainment among the intermediate-ability
group.  The escape-route incentive that was driving over-investment disappears;
these children now invest at their unconstrained optimum, which is strictly below
$\hat{e}$.  A policymaker monitoring Gaokao preparation intensity as a reform
indicator would observe declining effort among the most motivated rural students
precisely when the reform is working.

\section{Model}

\subsection{Agents and Ability}

Two children, an urban child $C_{U}$ and a rural child $C_{R}$, share innate
ability $\theta\in[\theta_{L},\theta_{H}]$ with $\theta_{L}>0$.  Hukou
registration status $\tau\in\{U,R\}$ is fixed at birth and determines
labor-market access.  Each child independently chooses education $e\in[0,\bar{e}]$
to maximize the present discounted value of lifetime earnings net of education
costs, discounted at rate $r>0$.  Agents are risk-neutral.

\begin{assumption}[Equal Ability]
Both children possess ability $\theta$; the investment gap
$\Delta e^{*}(\theta):=e^{*}_{U}(\theta)-e^{*}_{R}(\theta)$ is analyzed
parametrically in $\theta$ to isolate the institutional effect of hukou from
any ability differences between the urban and rural populations.
\end{assumption}

\subsection{Wages and the Hukou Penalty}

The market wage function is $w(e,\theta)=A\theta^{\alpha}e^{\beta}$, where
$A>0$, $\alpha>0$, and $\beta\in(0,1)$.  The Cobb-Douglas specification implies
strictly diminishing marginal returns to education ($w_{ee}<0$) and strict
supermodularity ($w_{e\theta}>0$): higher ability raises the marginal return to
education at every level $e$.

\begin{assumption}[Multiplicative Penalty]
The rural wage in the absence of an escape route is $\delta\cdot w(e,\theta)$,
$\delta\in(0,1)$.  The penalty scales the entire return to education
proportionally.
\end{assumption}

The escape-route probability function is
\begin{equation}
  p(e)=\begin{cases}0, & e<\hat{e},\\p_{0}, & e\geq\hat{e},\end{cases}
  \label{eq:escape}
\end{equation}
where $\hat{e}\in(0,\bar{e})$ is the elite education threshold and
$p_{0}\in(0,1)$.  In the Chinese context, $\hat{e}$ corresponds to crossing the
Gaokao provincial key-line for admission to a nationally recognized university,
and $p_{0}$ captures the fraction of elite rural graduates who successfully access
the urban labor market.  The discontinuity at $\hat{e}$ reflects the binary
nature of university-tier assignment.

\subsection{Payoffs}

Let $C(e)$ denote education costs with $C'>0$ and $C''>0$.  The urban child's
payoff is
\begin{equation}
  V_{U}(e;\theta)=\frac{w(e,\theta)}{r}-C(e).
  \label{eq:Vu}
\end{equation}
The rural child faces a two-regime structure.  Below the threshold,
\begin{equation}
  V_{R}^{low}(e;\theta)=\frac{\delta\,w(e,\theta)}{r}-C(e),\qquad e<\hat{e}.
  \label{eq:Vlow}
\end{equation}
At or above the threshold, the expected payoff reflects the escape-route lottery:
\begin{equation}
  \boxed{V_{R}^{high}(e;\theta)=
    \frac{\bigl[\delta+(1-\delta)p_{0}\bigr]w(e,\theta)}{r}-C(e),\qquad
    e\geq\hat{e}.}
  \label{eq:Vhigh}
\end{equation}
The effective wage multiplier $\delta+(1-\delta)p_{0}=1-(1-\delta)(1-p_{0})$
lies strictly between $\delta$ and $1$: crossing the elite threshold improves
but does not fully equalize expected wages, because the escape succeeds only with
probability $p_{0}<1$.

\section{Equilibrium Characterization}

\begin{lemma}[Existence and Characterization]
\label{lem:exist}
Under the Cobb-Douglas wage and strictly convex cost:
\begin{itemize}
  \item[\textit{(i)}] The urban child's optimum $e^{*}_{U}(\theta)$ uniquely
    solves
    \begin{equation}
      \frac{\beta A\theta^{\alpha}(e^{*}_{U})^{\beta-1}}{r}=C'(e^{*}_{U}).
      \label{eq:focU}
    \end{equation}
    Moreover, $e^{*}_{U}(\theta)$ is strictly increasing in $\theta$.
  \item[\textit{(ii)}] The rural child's global optimum satisfies
    \[
      e^{*}_{R}(\theta)\in
      \arg\max\Bigl\{\max_{e<\hat{e}}V_{R}^{low}(e;\theta),\;
      \max_{e\geq\hat{e}}V_{R}^{high}(e;\theta)\Bigr\}.
    \]
    The low-regime maximum is the unique interior point $e^{*}_{int}(\theta)$
    solving $\delta\,w_{e}(e^{*}_{int};\theta)/r=C'(e^{*}_{int})$.
    The high-regime maximum is at the interior point $e^{*}_{high}(\theta)$
    solving $[\delta+(1-\delta)p_{0}]\,w_{e}(e^{*}_{high};\theta)/r=
    C'(e^{*}_{high})$ when $e^{*}_{high}(\theta)\geq\hat{e}$, and at the corner
    $\hat{e}$ otherwise.  Since the effective multiplier
    $\delta+(1-\delta)p_{0}<1$, we have $e^{*}_{high}(\theta)<e^{*}_{U}(\theta)$
    for all $\theta$.  Recalling that $e^{*}_{U}(\bar{\theta})=\hat{e}$, the
    high-regime corner $\hat{e}$ is optimal within $[\hat{e},\bar{e}]$ for all
    $\theta\leq\bar{\theta}$, while for $\theta>\bar{\theta}$ the interior
    solution $e^{*}_{high}(\theta)>\hat{e}$ is feasible.
\end{itemize}
\end{lemma}

The urban child's problem is globally concave---$V_{U}$ is strictly concave in
$e$ because $w_{ee}<0$ and $C''>0$---so a unique interior solution exists.
Monotonicity of $e^{*}_{U}$ in $\theta$ follows from the supermodularity
condition $w_{e\theta}>0$: higher ability shifts the marginal-return schedule
upward, raising the crossing point with $C'(e)$.  The rural child cannot apply
a single first-order condition because $p(e)$ is discontinuous at $\hat{e}$.
The correct solution procedure is regime comparison: evaluate the best investment
in each regime and choose the higher-payoff option.  The over-investment result
(Proposition~\ref{prop:over}) focuses on the intermediate-ability range
$\theta\in(\tilde{\theta},\bar{\theta})$, where the high-regime maximum is at
the corner $\hat{e}$ and $\hat{e}>e^{*}_{U}(\theta)$.

\section{Under-Investment and Over-Investment}

\subsection{The Under-Investment Benchmark}

When the rural child's optimum lies below the escape threshold, the proportional
penalty delivers the \citet{Arrow1973} result in the human-capital setting.

\begin{proposition}[Under-Investment Benchmark]
\label{prop:under}
Suppose $e^{*}_{R}(\theta)<\hat{e}$.  Then
$e^{*}_{R}(\theta)<e^{*}_{U}(\theta)$ for all $\theta\in[\theta_{L},\theta_{H}]$.
The investment gap $\Delta e^{*}(\theta):=e^{*}_{U}(\theta)-e^{*}_{R}(\theta)$
is strictly positive and strictly decreasing in $\delta$, with
$\Delta e^{*}(\theta)\to 0$ as $\delta\to 1$.
\end{proposition}

The proof is a single-crossing argument.  The urban child's marginal return is
$w_{e}(e;\theta)/r$; the rural child's is $\delta\,w_{e}(e;\theta)/r<w_{e}/r$ at
every $e$.  Since $C'(e)$ is strictly increasing and identical across children,
the intersection of marginal benefit and marginal cost occurs at a strictly lower
$e$ for the rural child.  The mechanism is purely the return channel: the
multiplicative penalty scales down the payoff to every unit of education,
reducing investment.

\begin{remark}[Scope of Proposition~\ref{prop:under}]
\textit{The under-investment result depends critically on the multiplicative
penalty specification.  Under an additive penalty $W_{R}=w(e,\theta)-d$, the
rural child's marginal return is $w_{e}(e;\theta)/r$---identical to the urban
child's---so $e^{*}_{R}(\theta)=e^{*}_{U}(\theta)$ and the investment gap is
zero.  The multiplicative form is the mechanism through which the institutional
penalty distorts the education investment margin.}
\end{remark}

\subsection{Escape-Route Over-Investment}

The central result of this paper is that the direction of investment distortion
reverses for intermediate-ability rural children.  Define two ability thresholds:
\begin{align}
  \bar{\theta}&: \quad e^{*}_{U}(\bar{\theta})=\hat{e},
    \label{eq:thetabar}\\
  \tilde{\theta}&: \quad
    V_{R}^{low}\!\bigl(e^{*}_{int}(\tilde{\theta});\tilde{\theta}\bigr)=
    V_{R}^{high}\!\bigl(\hat{e};\tilde{\theta}\bigr).
    \label{eq:thetatilde}
\end{align}
$\bar{\theta}$ is the ability level at which the urban child's unconstrained
optimum equals $\hat{e}$; for $\theta<\bar{\theta}$, the urban child would invest
below $\hat{e}$ if unconstrained.  $\tilde{\theta}$ is the ability level at which
the rural child is indifferent between the low-regime interior solution and
investing exactly at $\hat{e}$.

\begin{proposition}[Escape-Route Over-Investment]
\label{prop:over}
If $p_{0}\geq p^{*}_{0}(\delta,\hat{e},\bar{\theta})$ (defined in the proof),
then $\tilde{\theta}<\bar{\theta}$ and
\[
  e^{*}_{R}(\theta)=\hat{e}>e^{*}_{U}(\theta)
  \quad\text{for all }\theta\in(\tilde{\theta},\bar{\theta}).
\]
A rural child of intermediate ability invests strictly more in education than an
equally-able urban child.
\end{proposition}

The proof has three steps.  \textit{Step~1.}  $\bar{\theta}$ is well-defined:
$e^{*}_{U}(\theta)$ is continuous and strictly increasing in $\theta$, and by
calibration $e^{*}_{U}(\theta_{L})<\hat{e}<e^{*}_{U}(\theta_{H})$, so the
intermediate value theorem delivers a unique $\bar{\theta}$.  \textit{Step~2.}
Define the payoff difference
$F(\theta):=V_{R}^{high}(\hat{e};\theta)-V_{R}^{low}(e^{*}_{int}(\theta);\theta)$.
At $\theta=\bar{\theta}$, the escape premium is
$(1-\delta)p_{0}\cdot w(\hat{e},\bar{\theta})/r>0$.  There exists a threshold
$p^{*}_{0}$ such that for $p_{0}\geq p^{*}_{0}$, this premium exceeds the net
cost of over-investing from $e^{*}_{int}(\bar{\theta})$ to $\hat{e}$, giving
$F(\bar{\theta})>0$.  \textit{Step~3.}  $F(\theta)$ is strictly increasing in
$\theta$: by supermodularity ($w_{e\theta}>0$) and the envelope theorem applied
to the low-regime optimum, higher ability raises $w(\hat{e},\theta)$ by more than
it raises $w(e^{*}_{int}(\theta),\theta)$, because $\hat{e}>e^{*}_{int}$.  Since
$F(\theta_{L})<0$ (negligible escape premium at low ability) and
$F(\bar{\theta})>0$, the intermediate value theorem delivers a unique
$\tilde{\theta}\in(\theta_{L},\bar{\theta})$ satisfying $F(\tilde{\theta})=0$.
For $\theta\in(\tilde{\theta},\bar{\theta})$, $F(\theta)>0$ implies
$e^{*}_{R}(\theta)=\hat{e}$.  Since $e^{*}_{U}(\theta)<\hat{e}$ for
$\theta<\bar{\theta}$, over-investment follows.

The economic mechanism is an asymmetric prize structure.  Urban children gain
nothing from crossing $\hat{e}$: they already earn the full urban wage at any
education level.  Rural children gain an expected wage jump of
$(1-\delta)p_{0}\cdot w(\hat{e},\theta)/r$ upon crossing.  For intermediate-ability
rural children whose unconstrained optimum lies just below $\hat{e}$, this escape
premium outweighs the marginal cost of the additional investment required to reach
the threshold.  The result overturns the central prediction of standard
discrimination models---that institutional barriers uniformly suppress investment
by the affected group---by showing that a threshold-based escape route creates
over-investment incentives for those close enough to the threshold to find the
gamble rational.

\begin{remark}[Robustness to Penalty Specification]
\textit{Unlike Proposition~\ref{prop:under}, the over-investment result is robust
to the penalty specification.  Under an additive penalty $W_{R}=w-d$, the escape
premium at $\hat{e}$ becomes $d\cdot p_{0}/r>0$, which still creates an incentive
to invest at $\hat{e}$ for intermediate-ability rural children.  The
over-investment direction survives; the under-investment direction does not.}
\end{remark}

\section{Comparative Statics}

\begin{proposition}[Comparative Statics on the Investment Gap]
\label{prop:cs}
\begin{itemize}
  \item[\textit{(i)}] For $\theta\in[\theta_{L},\tilde{\theta})$:
    $\partial\Delta e^{*}(\theta)/\partial\delta<0$.
    Relaxing the hukou penalty reduces the under-investment gap.
  \item[\textit{(ii)}] $\partial(\bar{\theta}-\tilde{\theta})/\partial p_{0}>0$.
    A higher escape probability expands the ability range over which rural
    children over-invest.
  \item[\textit{(iii)}] The sign of $\partial\tilde{\theta}/\partial\delta$
    depends on a parameter condition.  By the implicit function theorem and the
    envelope theorem applied to $F(\theta):=V_{R}^{high}(\hat{e};\theta)-
    V_{R}^{low}(e^{*}_{int}(\theta);\theta)$:
    \[
      \operatorname{sign}\!\Bigl(\frac{\partial\tilde{\theta}}{\partial\delta}
      \Bigr)
      =-\operatorname{sign}\!\bigl[(1-p_{0})\hat{e}^{\beta}-(e^{*}_{int})^{\beta}
      \bigr].
    \]
    When $\hat{e}/e^{*}_{int}(\tilde{\theta})>(1-p_{0})^{-1/\beta}$:
    $\partial\tilde{\theta}/\partial\delta<0$, so a \emph{relaxation} of the
    hukou penalty (higher $\delta$) lowers $\tilde{\theta}$ and expands the
    over-investment range.  Under the complementary condition
    ($\hat{e}/e^{*}_{int}\leq(1-p_{0})^{-1/\beta}$): $\partial\tilde{\theta}/
    \partial\delta\geq 0$ and stricter enforcement widens the range.
    The aggregate effect of enforcement on the over-investment range is
    therefore parameter-dependent and not generically monotone.
\end{itemize}
\end{proposition}

Parts~(i) and~(iii) follow from the implicit function theorem applied to the
low-regime first-order condition and indifference condition~(\ref{eq:thetatilde}),
respectively.  Part~(ii) follows because $F(\theta)$ is strictly increasing in
$p_{0}$ at every $\theta$, which shifts $\tilde{\theta}$ downward while leaving
$\bar{\theta}$ unchanged ($\bar{\theta}$ depends only on $e^{*}_{U}$, which is
independent of $p_{0}$).

The sign result in part~(iii) has a clean intuition.  When $\hat{e}\gg
e^{*}_{int}$ (first condition), a marginal increase in $\delta$ raises the
high-regime base return $\delta w(\hat{e},\theta)$ by much more than it raises
the low-regime return $\delta w(e^{*}_{int},\theta)$, since $\hat{e}^{\beta}\gg
(e^{*}_{int})^{\beta}$.  Even though the escape premium $(1-\delta)p_{0}w$ falls,
the net payoff advantage of the high regime increases, drawing more rural children
toward over-investment.  When $\hat{e}$ is close to $e^{*}_{int}$ (second
condition), the escape premium effect dominates and stricter enforcement widens
the range.  Neither direction is empirically universal; calibration to Chinese
Gaokao data is required to determine which regime obtains.

\section{Welfare}

Let $e^{FB}(\theta)=e^{*}_{U}(\theta)$ denote the first-best investment level
and define the welfare loss from hukou as
$\Delta W(\theta):=V_{U}(e^{FB};\theta)-V_{R}(e^{*}_{R};\theta)$.

\begin{proposition}[Welfare Loss]
\label{prop:welfare}
For all $\theta\in[\theta_{L},\theta_{H}]$ and $\delta\in(0,1)$:
\begin{itemize}
  \item[\textit{(i)}] $\Delta W(\theta)>0$.
  \item[\textit{(ii)}] $\partial\Delta W(\theta)/\partial\delta<0$: relaxing the
    penalty reduces welfare loss.
  \item[\textit{(iii)}] $\Delta W(\theta)$ is non-monotone in $\theta$, attaining
    a local maximum near $\tilde{\theta}$.
\end{itemize}
\end{proposition}

The non-monotone profile in part~(iii) reflects two competing distortions.  In
the under-investment region $\theta<\tilde{\theta}$, the gap between $e^{FB}$ and
$e^{*}_{R}$ grows as $\theta$ approaches $\tilde{\theta}$ from below:
$e^{*}_{U}(\theta)$ rises toward $\hat{e}$ while $e^{*}_{int}(\theta)$ is pushed
downward by the penalty.  In the over-investment region $\theta\in(\tilde{\theta},
\bar{\theta})$, the rural child invests at $\hat{e}$ and bears the cost of
over-shooting with certainty, while receiving the escape premium only with
probability $p_{0}$.  The welfare loss peaks near $\tilde{\theta}$ where both
distortions are simultaneously active.

A natural prior predicts that welfare losses should be largest for high-ability
rural children, who have the most to gain from equal treatment.
Proposition~\ref{prop:welfare}(iii) qualifies this: above $\bar{\theta}$, rural
children invest in the high regime at $e^{*}_{high}(\theta)>\hat{e}$, but the
effective multiplier $\delta+(1-\delta)p_{0}<1$ means their expected wage still
falls short of the first-best; $\Delta W(\theta)>0$ for all $\theta$.  What the
non-monotonicity says is that the welfare loss is particularly large near
$\tilde{\theta}$, where two distortions coincide: the rural child over-shoots to
$\hat{e}$ (incurring excess investment cost) while still receiving the penalized
wage with probability $1-p_{0}$.  Whether the welfare loss is monotonically
smaller above $\bar{\theta}$ than near $\tilde{\theta}$ depends on parameters and
is not established in general; the proposition asserts only that $\Delta W(\theta)$
attains a local maximum near $\tilde{\theta}$, not a global ordering across the
full ability distribution.

\begin{proposition}[Unintended Consequence of Hukou Reform]
\label{prop:reform}
Let $\delta^{*}\in(0,1)$ be the pre-reform penalty.  A reform setting $\delta'=1$:
\begin{itemize}
  \item[\textit{(i)}] Eliminates all welfare loss: $\Delta W(\theta;\delta')=0$
    for all $\theta$.
  \item[\textit{(ii)}] Increases investment for
    $\theta\in[\theta_{L},\tilde{\theta}(\delta^{*}))$:
    $e^{*}_{R}(\theta;\delta')=e^{*}_{U}(\theta)>e^{*}_{R}(\theta;\delta^{*})$.
  \item[\textit{(iii)}] \textit{Decreases} investment for
    $\theta\in(\tilde{\theta}(\delta^{*}),\bar{\theta})$:
    $e^{*}_{R}(\theta;\delta')=e^{*}_{U}(\theta)<\hat{e}=
    e^{*}_{R}(\theta;\delta^{*})$.
\end{itemize}
\end{proposition}

The proof is direct.  Post-reform ($\delta=1$), both children solve the same
problem and arrive at the first-best: $e^{*}_{R}(\theta;1)=e^{*}_{U}(\theta)$.
Welfare is equalized by construction.  The investment comparison in~(iii) follows
because the pre-reform rural child in $(\tilde{\theta},\bar{\theta})$ was
investing at $\hat{e}>e^{*}_{U}(\theta)$, while the post-reform rural child
invests at $e^{*}_{U}(\theta)$.

The paradox arises because the escape-route over-investment was a distortion, not
genuine educational ambition.  Once the wage penalty disappears, intermediate-ability
rural children optimally reduce education to the unconstrained level: they earn the
full wage at lower cost, strictly improving welfare.  A policymaker who interprets
reduced Gaokao preparation intensity among rural near-threshold students as evidence
that reform has backfired is drawing the opposite conclusion from the correct one:
the observed investment decline is the efficiency gain from removing the distortion.

\section{Boundary Cases}

\begin{proposition}[Limiting Behavior]
\label{prop:boundary}
\begin{itemize}
  \item[\textit{(i)}] As $\delta\to 1$: $e^{*}_{R}(\theta)\to e^{*}_{U}(\theta)$
    and $\Delta W(\theta)\to 0$ for all $\theta$.
  \item[\textit{(ii)}] As $p_{0}\to 0$: the over-investment range collapses,
    $\bar{\theta}-\tilde{\theta}\to 0$.  For $\theta\leq\bar{\theta}$,
    $e^{*}_{R}(\theta)\to e^{*}_{int}(\theta)$ (the low-regime interior).
    For $\theta>\bar{\theta}$, the high-regime interior $e^{*}_{high}(\theta)\to
    e^{*}_{int}(\theta)$ as well (since the effective multiplier
    $\delta+(1-\delta)p_{0}\to\delta$), representing a \emph{contraction} of
    investment for these high-ability children.  The overall limit is pure
    under-investment relative to first-best for all $\theta$, consistent with
    the \citet{Arrow1973} benchmark.
  \item[\textit{(iii)}] As $\delta\to 0$: $e^{*}_{int}(\theta)\to 0$ for all
    $\theta$; the escape route becomes the only investment incentive for rural
    children.  Whether rural children actually cross $\hat{e}$ depends on whether
    the escape premium $p_{0}\,w(\hat{e},\theta)/r$ exceeds $C(\hat{e})$; for
    very low $\delta$ and low $\theta$, neither regime may generate positive
    investment.
\end{itemize}
\end{proposition}

The limit in~(ii) is the most informative.  As $p_{0}\to 0$, the payoff jump at
$\hat{e}$ vanishes; the rural child always prefers the low-regime interior,
yielding the Arrow~(1973) prediction of proportional under-investment at every
ability level.  For high-ability children ($\theta>\bar{\theta}$), convergence
to the Arrow benchmark involves a fall in investment from the high-regime interior
$e^{*}_{high}>\hat{e}$ back to $e^{*}_{int}<\hat{e}$, not merely a collapse of
the intermediate over-investment range.  The novel predictions of the present
paper---escape-route over-investment (Proposition~\ref{prop:over}) and the reform
paradox (Proposition~\ref{prop:reform})---require both $\delta<1$ and $p_{0}>0$;
neither feature alone is sufficient.

\section{Testable Predictions}

The model generates the following empirically testable predictions.

\textbf{P1.~Non-monotone investment gap.}  Among rural and urban students of
equal measured ability, the investment gap $\Delta e^{*}(\theta)$ is positive
(rural under-invest) at low ability, negative (rural over-invest) at intermediate
ability near the Gaokao key-line, and approaches zero at high ability.
Cross-sectional variation in Gaokao preparation intensity by ability and hukou
status provides the most direct test.

\textbf{P2.~Over-investment expands under stricter enforcement.}  When the elite
threshold substantially exceeds rural students' unconstrained optimum
(Proposition~\ref{prop:cs}(iii)), stricter regional hukou enforcement should
increase preparation intensity among near-threshold rural students.  Regional
variation in enforcement strictness is the natural testing ground.

\textbf{P3.~Post-reform investment decline among near-threshold students.}  In
regions where hukou reform has reduced the wage penalty, Gaokao preparation
intensity among intermediate-ability rural students should fall even as wage and
welfare outcomes improve.  Isolating this prediction from general discouragement
effects requires joint measurement of wage outcomes and investment.

\textbf{P4.~Welfare loss concentrated at intermediate ability.}  Welfare measures
---earnings gaps, consumption shortfalls, subjective well-being differentials
relative to urban peers---should be largest for rural students near the elite
admission threshold, not for the most- or least-able.

\begin{thebibliography}{}

\bibitem[Arrow(1973)]{Arrow1973}
Arrow, K. J. 1973. ``The Theory of Discrimination.''
In \textit{Discrimination in Labor Markets},
edited by O. Ashenfelter and A. Rees, 3--33.
Princeton University Press.

\bibitem[Becker(1964)]{Becker1964}
Becker, G. S. 1964.
\textit{Human Capital: A Theoretical and Empirical Analysis,
with Special Reference to Education}.
National Bureau of Economic Research.

\bibitem[Galor and Zeira(1993)]{GalorZeira1993}
Galor, O., and J. Zeira. 1993.
``Income Distribution and Macroeconomics.''
\textit{Review of Economic Studies}, 60(1): 35--52.

\end{thebibliography}

\end{document}
